\documentclass{article}
\usepackage{/globals/anki}
\ankiprefix{beratung}
\ankitopic{Beratung}
\ankishow
\begin{document}
\section{Beratung}
BeraterInnen arbeiten als Single Point of Contact für Kunden und steuern mit projektübergreifendem Wissen Projekte, bereiten diese vor und nach, führen sie durch und halten zeitliche und finanzielle Vorgaben im Blick.
\ankinote{aufgaben}{Aufgaben}{Single Point of Contact, Wissen, Projekte, Vorgaben}

\subsection{Ablauf eines Projektes}
Ein Projekt besteht aus der Anreise zum Kunden, der Besprechung dessen Anforderungen, der Entwicklung von passenden Konzepten, dem Gespräch mit dem Kunden, der Digitalisierung der Prozesse und der nachfolgenden Schulung der MitarbeiterInnen für die Software und der Betreuung der Applikation.
\ankinote{prozess}{Prozess}{Anforderungen, Konzepte, Gespräch, Digitalisierung, Schulung, Betreuung}

\subsection{Releasemanagement}
Beim Releasemanagement wird die Produktivsetzung, die Veröffentlichung eines Produktes, geplant, koordiniert und mit durchgeführt. Dabei wird sichergestellt dass alle Vorgaben eingehalten werden. Zudem muss oft bei Fragen oder Problemen unterstützt werden.
\ankinote{release}{Releasemanagement}{Produktivsetzung planen, durchführen; Vorgaben; Unterstützung}



\end{document}